\clearpage
\VerdierPhysicalPageAnchor{0231}%
\begingroup\small\setlength{\parskip}{0pt}\setlength{\emergencystretch}{1.5em}
\setlength{\abovedisplayskip}{4pt plus 2pt minus 2pt}%
\setlength{\belowdisplayskip}{4pt plus 2pt minus 2pt}%

\VerdierSemanticAnchor{ch4d-4-1-7-proof}
Suppose first that the proposition has been proved when $X$ is a double
complex and the integer $m=1$, and let us show how one can reduce to this
case. Let $\varphi\colon[n]\longrightarrow[2]$ be the map:

\VerdierSemanticAnchor{ch4d-4-1-7-1}
\begin{equation}
\left\{
\begin{aligned}
\varphi(m)&=1,\\
\varphi(i)&=2,\qquad i\in[n],\quad i\neq m.
\end{aligned}
\right.
\tag{4.1.7.1}
\end{equation}

Set $\int_\varphi X=Z$ (Chapter~I, 2.2.2). One may then carry out on the
double complex $Z$ the constructions described in (4.1.6.$i$),
$1\leq i\leq4$. One obtains double complexes

\VerdierSemanticAnchor{ch4d-4-1-7-2}
\begin{equation}
{}^{1}\!Z(p,q),\qquad
(p,q)\in\mathcal{F}\ell_1(\widetilde{\mathbb Z}),
\tag{4.1.7.2}
\end{equation}

and morphisms of double complexes

\VerdierSemanticAnchor{ch4d-4-1-7-3}
\begin{equation}
\begin{aligned}
f_Z\bigl((p,q),(p',q')\bigr)&\colon
{}^{1}\!Z(p,q)\longrightarrow{}^{1}\!Z(p',q'),\\
g_Z(p,q,r)&\colon
{}^{1}\!Z(q,r)\longrightarrow T_1\bigl({}^{1}\!Z(p,q)\bigr).
\end{aligned}
\tag{4.1.7.3}
\end{equation}

One then sees that

\VerdierSemanticAnchor{ch4d-4-1-7-4}
\begin{equation}
\left\{
\begin{aligned}
\int_\varphi{}^{m}\!X(p,q)
  &={}^{1}\!Z(p,q),
  &&(p,q)\in\mathcal{F}\ell_1(\widetilde{\mathbb Z}),\\
\int_\varphi f\bigl((p,q),(p',q')\bigr)
  &=f_Z\bigl((p,q),(p',q')\bigr),
\end{aligned}
\right.
\tag{4.1.7.4}
\end{equation}

(by direct verification), and moreover the morphism $g_Z(p,q,r)$ is the
composite morphism

\VerdierSemanticAnchor{ch4d-4-1-7-5}
\begin{equation}
\int_\varphi{}^{m}\!X(q,r)
\xrightarrow{\int_\varphi g(p,q,r)}
\int_\varphi T_m\bigl({}^{m}\!X(p,q)\bigr)
\xrightarrow{(**)}
T_1\left(\int_\varphi{}^{m}\!X(p,q)\right)
\tag{4.1.7.5}
\end{equation}

where $(**)$ is the isomorphism expressing commutation of the functor
$\int_\varphi$ with translations (Chapter~I, 2.3.2.1) (to verify this one
uses the explicit definition of the differentials of $Z$ (Chapter~I,
2.2.2.5) and of the commutation isomorphism (Chapter~I, 2.3.2.1). Now
construct the spectral object defined by $Z$ filtered by its first degree.
It follows from the equality $\int_2\int_\varphi=\int_n$, from (4.1.7.4),
(4.1.7.5), and from (Chapter~I, 2.3.3) that
$\bigl({}^{m}\!Y(p,q),{}^{m}\!d(p,q,r)\bigr)$ is equal to this spectral
object.
\endgroup
